% DropRebound_ContactLine.tex
% Companion modeling derivation: contact-line dynamics (finite equilibrium
% contact angle + linear contact-line mobility) for the kinematic-match
% solid-substrate rebound model.
%
% This document extends ONLY the modeling portion of
%   DropRebound_JFM.tex  ("Drop rebound at low Weber number")
% specifically Sec. 3.1 (Kinematic match model formulation), Appendix A
% (Defining contact) and Appendix B (Numerical approximations).
%
% STATUS: gold-standard derivation, to be certified by independent adversarial
% verification BEFORE any code is written. Notation deliberately matches the
% parent paper.

\documentclass{jfm}
\usepackage{natbib}
\usepackage{graphicx}
\usepackage{amsmath,mathtools}
\usepackage{amsfonts}
\usepackage{amssymb}
\usepackage{upgreek}
\usepackage{mathrsfs}
\usepackage{hyperref}

\newcommand{\be}{\begin{equation}}
\newcommand{\ee}{\end{equation}}
\newcommand{\We}{W\kern-0.15em e\,}
\newcommand{\Bo}{B\kern-0.05em o\,}
\newcommand{\Oh}{O\kern-0.05em h\,}
\newcommand{\Ca}{C\kern-0.05em a\,}

\newtheorem{lemma}{Lemma}
\newtheorem{proposition}{Proposition}
\newtheorem{corollary}{Corollary}
\newtheorem{remark}{Remark}

\title{Contact-line dynamics in the kinematic-match model of drop rebound:
finite contact angle and linear contact-line friction}

\shorttitle{Contact-line dynamics in kinematic-match rebound}
\shortauthor{E.\,A. Ag\"uero and others}

\author{Elvis A. Ag\"uero\aff{1} \and others}
\affiliation{\aff{1}School of Engineering, Brown University, Providence, RI 02912, USA}

\begin{document}
\maketitle

\begin{abstract}
The kinematic-match (KM) reduced-order model of \citet{gabbard2025} (hereafter GA)
predicts drop rebound from a \emph{perfectly non-wetting} substrate with no fitting
parameters. Its non-wetting character is encoded in a single boundary condition: the
free surface departs the substrate tangentially at the contact line. We relax this
assumption from first principles, admitting substrates of finite equilibrium contact
angle $\theta_e$ and a finite contact-line mobility. We frame the problem in a
metriplectic-style (\textsc{generic}-flavoured) form: the reversible capillary--inertial mode oscillator
of GA, a unilateral non-penetration constraint whose Lagrange multiplier is the contact
pressure, and a new \emph{irreversible} dissipation localised at the moving contact
line. Two results follow. First, the tangency condition is replaced by a dynamic
contact-angle condition; in the quasi-static (perfectly mobile) limit this reduces to
imposing the apparent contact angle $\theta_d=\theta_e$, and for $\theta_e\to\pi$ it
recovers GA's tangency exactly. Second, the uncompensated Young force is radial and
cancels by axisymmetry, so \emph{the centre-of-mass momentum equation of GA is
unchanged}; the entire wetting effect enters through (i) the contact-line boundary
condition and (ii) a contact-line dissipation channel governed by a linear (Blake/de
Gennes) friction law $\xi\,\dot r_c=\sigma(\cos\theta_e-\cos\theta_d)$. The extended
model has exactly one new dynamical degree of freedom (the contact-line radius $r_c$)
and two physical parameters, the measurable $\theta_e$ and a single friction
coefficient $\xi$; it reduces to GA in the limit $(\theta_e\to\pi,\ \xi\to0)$. We give
the discrete linear system, showing the extension changes only the contact-set
\emph{selection criterion} (quasi-static) plus a rate constraint (dynamic), preserving
the linearity that underlies GA's numerical robustness.
\end{abstract}

\section{Introduction}\label{sec:intro}

A drop striking a solid or liquid surface may rebound without ever coalescing, a
phenomenon central to superhydrophobic surface design, spray cooling, inkjet printing
and self-cleaning materials \citep{Richard2002,SnoeijerAndreotti2013}. When rebound is
mediated by a thin gas film the impact is effectively \emph{non-wetting}: no
three-phase contact line forms, and the drop's contact time and coefficient of
restitution are set by the interplay of inertia, capillarity and viscosity
\citep{Richard2002,Okumura2003,MolacekBush2013}. Reduced-order ``kinematic-match'' (KM)
models exploit this by evolving the drop's spherical-harmonic shape modes subject to a
unilateral, tangential contact with the substrate, and reproduce experiment across a
wide range of Weber number at a small fraction of the cost of direct numerical
simulation \citep{GaleanoRios2017,Alventosa2023,gabbard2025}.

The defining assumption of these models is \emph{perfect} non-wetting: the free surface
meets the substrate tangentially, so the equilibrium contact angle is $\theta_e=\pi$,
there is no adhesion, and the contact line carries no localised dissipation. Real
substrates --- including the hydrophobic coatings used to suppress, but not eliminate,
wetting --- depart from this idealisation. A finite equilibrium contact angle introduces
an uncompensated Young force at the contact line \citep{deGennes1985}, its motion
dissipates energy through microscopic processes captured at leading order by a linear
(molecular-kinetic / hydrodynamic) friction law
\citep{Blake2006,Cox1986,Voinov1976,SnoeijerAndreotti2013}, and the advancing and
receding branches generically differ, producing contact-angle hysteresis
\citep{JoannyDeGennes1984}. None of these ingredients is present in the KM framework.

Here we extend the KM model to substrates of finite wettability from first principles,
retaining the linear structure responsible for its numerical robustness. We organise the
dynamics in a metriplectic (\textsc{generic}-flavoured) form
\citep{GrmelaOttinger1997}: the reversible capillary--inertial mode oscillator of
\citet{gabbard2025} (hereafter GA), a unilateral non-penetration constraint whose
Lagrange multiplier is the contact pressure, and an irreversible dissipation localised
at the contact line. The tangency condition is replaced by a dynamic contact-angle
condition closed by a Blake/de Gennes mobility law; the perfectly non-wetting model is
recovered as the limit $\theta_e\to\pi$ with infinite mobility. We show that the
extension is a change of \emph{constraint} rather than an added body force, so the
centre-of-mass equation and the discrete linear system of GA are structurally preserved.

\section{The perfectly non-wetting model}\label{sec:background}

We adopt the notation and non-dimensionalisation of GA Sec.~3.1: lengths by the
undeformed radius $R$, mass by $\rho R^3$, time by the inertio-capillary time
$t_\sigma=(\rho R^3/\sigma)^{1/2}$, so that $\sigma\equiv1$ in dimensionless units and
gravity and viscosity are carried by $\Bo=\rho g R^2/\sigma$ and
$\Oh=\mu/\sqrt{\rho\sigma R}$. The axisymmetric interface is
\be\label{eq:shape}
\zeta(\theta,t)=1+\eta(\theta,t),\qquad
\eta=\sum_{l=2}^{\infty}\mathscr{A}_l(t)\,P_l(\cos\theta),
\ee
with $P_l$ the Legendre polynomials and $\theta\in[0,\pi]$ the polar angle measured from
the downward vertical, so that $\theta=0$ is the pole facing the substrate. The contact
pressure is $p(\theta,t)=\sum_{l\ge0}\mathscr{B}_l(t)P_l(\cos\theta)$. With
$\mathscr{U}_l=\dot{\mathscr A}_l$ the modal surface velocity, $h$ the centre-of-mass
(COM) height above the substrate and $v=\dot h$, the linearised dynamics of GA are
\begin{align}
\dot{\mathscr A}_l &= \mathscr U_l, \label{eq:kin}\\
\dot{\mathscr U}_l &= -\,l(l+2)(l-1)\,\mathscr A_l
                      -2(2l+1)(l-1)\,\Oh\,\mathscr U_l
                      -\,l\,\mathscr B_l, \label{eq:mode}\\
\dot h &= v, \label{eq:hdot}\\
\dot v &= -\Bo + \mathscr B_1. \label{eq:vdot}
\end{align}
Equations \eqref{eq:kin}--\eqref{eq:vdot} hold for $l=2,\dots,\infty$. The COM force
\eqref{eq:vdot} is the net vertical reaction of the contact pressure (its $l=1$ moment)
plus gravity; we return to this in \S\ref{sec:noforce}.

\medskip\noindent\textbf{Cartesian geometry.} A surface point at polar angle $\theta$
sits at cylindrical radius and height
\be\label{eq:cart}
x(\theta)=\zeta(\theta)\sin\theta,\qquad
z(\theta)=h-\zeta(\theta)\cos\theta,
\ee
so the pole $\theta=0$ is at height $z=h-\zeta(0)$ and the substrate is $z=0$. The
signed gap to the substrate is $z(\theta)$; contact requires $z=0$ and impenetrability
requires $z\ge 0$.

\medskip\noindent\textbf{The kinematic match.} Where the substrate presses the drop a
contact cap $\theta\le\theta_c$ forms on $z=0$:
\begin{align}
\sum_{l\ge2}\mathscr A_l P_l(\cos\theta) &= \frac{h}{\cos\theta}-1,
   &&\theta\le\theta_c, \label{eq:contact}\\
\sum_{l\ge2}\mathscr A_l P_l(\cos\theta) &< \frac{h}{\cos\theta}-1,
   &&\theta>\theta_c, \label{eq:noncontact}\\
\sum_{l\ge0}\mathscr B_l P_l(\cos\theta) &= 0,
   &&\theta>\theta_c, \label{eq:pzero}
\end{align}
where \eqref{eq:noncontact} is the strict form of $z>0$ off contact (equivalently
$\zeta\cos\theta<h$).\footnote{The parent paper GA writes the non-contact inequality as
$\zeta\cos\theta>h$ (their eqns for \eqref{eq:noncontact} and the impenetrability
predicate $e^1_q$); evaluated at the drop apex $\theta=\pi$, where $z=h+\zeta>0$
unambiguously, that form reads $-\zeta>h$, which is false. The correct sense is
$\zeta\cos\theta<h$ as used here; GA's printed inequality is a sign typo (the two
formulations agree, as they must, on the contact \emph{equality} $\zeta\cos\theta=h$).}
Finally GA impose \emph{tangential departure} at the contact line,
\be\label{eq:tangency}
\left.\frac{\partial}{\partial\theta}\Big[\zeta(\theta)\cos\theta\Big]\right|_{\theta=\theta_c}=0,
\ee
which, since $\partial_\theta(\zeta\cos\theta)=-\,\partial_\theta z$, states
$\partial_\theta z|_{\theta_c}=0$: the gap has a smooth (tangential) zero. This is the
sole encoding of perfect non-wetting; \S\ref{sec:angle} replaces it.

\section{The apparent contact angle}\label{sec:angle}

Define the \emph{apparent (dynamic) contact angle} $\theta_d$ as the angle, measured
through the liquid, between the substrate and the liquid--gas interface at the contact
line. Using \eqref{eq:cart}, the interface tangent in the direction of increasing
$\theta$ (i.e.\ away from the wetted pole, up the free surface) is
$\mathbf T=(x_\theta,z_\theta)$ with, writing
$\zeta_\theta\equiv\partial_\theta\zeta=-\sin\theta\sum_{l\ge2}\mathscr A_l P_l'(\cos\theta)$,
\be\label{eq:tangent}
x_\theta=\zeta_\theta\sin\theta+\zeta\cos\theta,\qquad
z_\theta=\zeta\sin\theta-\zeta_\theta\cos\theta .
\ee
Let $\alpha=\operatorname{atan2}(z_\theta,x_\theta)$ be the inclination of $\mathbf T$
above the substrate ($+x$). The liquid occupies the interior (axis) side, so the contact
angle through the liquid is the supplement,
\be\label{eq:thetad}
\boxed{\;\theta_d=\pi-\alpha
=\pi-\operatorname{atan2}\!\big(\zeta\sin\theta_c-\zeta_\theta\cos\theta_c,\;
\zeta_\theta\sin\theta_c+\zeta\cos\theta_c\big).\;}
\ee

\begin{lemma}[Undeformed reduction]\label{lem:sphere}
For an undeformed sphere ($\mathscr A_l\equiv0$, so $\zeta=1$, $\zeta_\theta=0$),
$\theta_d=\pi-\theta_c$. In particular $\theta_c\to0$ (point contact) gives
$\theta_d\to\pi$ and $\theta_c\to\pi/2$ (hemisphere) gives $\theta_d\to\pi/2$.
\end{lemma}
\begin{proof}
With $\zeta=1,\zeta_\theta=0$, \eqref{eq:tangent} gives $(x_\theta,z_\theta)=(\cos\theta_c,\sin\theta_c)$,
so $\alpha=\theta_c$ and $\theta_d=\pi-\theta_c$.
\end{proof}

\begin{proposition}[Tangency is the $\theta_d=\pi$ case]\label{prop:tangency}
Assume the contact line lies on the ascending flank of the profile, $x_\theta(\theta_c)>0$
(the cylindrical radius increases outward there). Then the GA tangency condition
\eqref{eq:tangency} holds iff $\theta_d=\pi$.
\end{proposition}
\begin{proof}
$\partial_\theta z=z_\theta=\zeta\sin\theta_c-\zeta_\theta\cos\theta_c$ is the first
argument of $\operatorname{atan2}$ in \eqref{eq:thetad}. Under the hypothesis
$x_\theta>0$ the second argument is positive, so
$z_\theta=0\iff\alpha=0\iff\theta_d=\pi$. By \eqref{eq:tangency},
$\partial_\theta z|_{\theta_c}=0$, hence $\theta_d=\pi$, and conversely.
\end{proof}

\begin{remark}[Validity of the $x_\theta>0$ hypothesis]\label{rem:xtheta}
$x_\theta=\partial_\theta(\zeta\sin\theta)$ is the slope of the cylindrical radius and
vanishes at the profile's widest point: for the undeformed sphere $x_\theta=\cos\theta_c$,
which is positive for $\theta_c<\pi/2$, zero at $\theta_c=\pi/2$, and negative beyond.
The hypothesis therefore holds throughout the low-$\We$, small-contact regime
($\theta_c<\pi/2$) that is our target, but can fail near maximal spreading or for
pancake/overhanging shapes, where $\alpha=\operatorname{atan2}(z_\theta,x_\theta)$
becomes ill-conditioned. The discrete residual of \S\ref{sec:discrete} is deliberately
written in a form \eqref{eq:disc_angle} that avoids dividing by $x_\theta$ and so
remains well defined there.
\end{remark}

Thus GA is exactly the special case $\theta_d\equiv\pi$. A substrate of finite
wettability instead selects a finite apparent angle, determined dynamically below.

\section{Energy structure and the wetting energy}\label{sec:energy}

Organise the model in a metriplectic-style (\textsc{generic}-flavoured) form: reversible
kinetic$+$surface energy, a unilateral constraint, and irreversible dissipation. We do
not verify the formal \textsc{generic} degeneracy conditions; the structure is used as an
organising principle and is \emph{backed} by the closed energy budget of
\S\ref{sec:budget}. The total energy is
\be\label{eq:energy}
E=\underbrace{\tfrac12\!\int_{\text{drop}}\!|\mathbf u|^2\,\mathrm dV}_{E_{\text{kin}}}
 +\underbrace{\Bo\,h}_{E_{\text{grav}}}
 +\underbrace{A_{lg}}_{\text{liquid--gas}}
 +\underbrace{(\gamma_{sl}-\gamma_{sv})\,A_{w}}_{\text{wetting}},
\ee
in dimensionless units ($\sigma=1$; $\gamma_{sl},\gamma_{sv}$ the solid--liquid and
solid--vapour tensions in units of $\sigma$). Young's relation
$\gamma_{sv}-\gamma_{sl}=\cos\theta_e$ turns the wetting term into
\be\label{eq:Ewet}
E_{\text{wet}}=-\cos\theta_e\,A_w=-\pi\cos\theta_e\,r_c^{2},\qquad
r_c\equiv x(\theta_c)=\zeta(\theta_c)\sin\theta_c,
\ee
$r_c$ being the contact-line radius and $A_w=\pi r_c^2$ the wetted area to leading order.
The liquid--gas area, expanded to quadratic order in the $\mathscr A_l$, reproduces the
Lamb capillary spectrum and hence the restoring term $-l(l+2)(l-1)\mathscr A_l$ already
present in \eqref{eq:mode} \citep{Lamb1924}; no new bulk term arises. The genuinely new
ingredient is \eqref{eq:Ewet}.

\begin{proposition}[Reduction of the wetting energy]\label{prop:Ewet}
As $\theta_e\to\pi$, $\cos\theta_e\to-1$ and $E_{\text{wet}}\to\pi r_c^2$; its variations
drive the contact line only through \eqref{eq:mobility} below, whose driving force
$\propto(\cos\theta_e-\cos\theta_d)$ vanishes when additionally $\theta_d\to\pi$
(Proposition~\ref{prop:tangency}). Hence at $(\theta_e\to\pi)$ with a tangential contact
line the wetting contribution is dynamically inert, consistent with GA.
\end{proposition}

\section{The centre-of-mass equation retains the GA form}\label{sec:noforce}

A central and non-obvious claim: the COM equation keeps the GA \emph{form}
$\dot v=-\Bo+\mathscr B_1$ even for finite wettability --- not because no capillary
force acts at the contact line, but because that force is reacted by the constraint
pressure and thereby absorbed into the pressure moment $\mathscr B_1$.

\begin{proposition}[COM equation form is preserved]\label{prop:noforce}
The vertical force balance of the drop's centre of mass retains the GA form
$\dot v=-\Bo+\mathscr B_1$, where $\mathscr B_1$ is the $l=1$ moment of the contact
pressure. In particular the contact-line capillary reaction does not appear as a
\emph{separate} term.
\end{proposition}
\begin{proof}
Take the drop (liquid plus its bounding liquid--gas interface) as the system. The
external agents are gravity, the uniform vapour pressure (absorbed into the reference),
and the solid. At the contact line the liquid--gas tension exerts a line force of
magnitude $\sigma$ per unit length directed along the interface tangent; its
\emph{radial} (in-plane) component is the uncompensated Young force
$\sigma(\cos\theta_e-\cos\theta_d)$ per length, which is horizontal and cancels by
axisymmetry, while its \emph{vertical} component is $\sigma\sin\theta_d$ per length, a
net upward line load
\be\label{eq:linelift}
F_{\text{cl}}^{z}=2\pi r_c\,\sigma\sin\theta_d,
\ee
nonzero for any $\theta_d\neq\pi$. This vertical pull is not balanced by
$\gamma_{sl},\gamma_{sv}$ (which are horizontal); it is balanced by the substrate
\emph{normal reaction}, i.e.\ by the contact pressure $p$, which is the Lagrange
multiplier enforcing $z=0$ on the wetted cap. In the reduced model the net vertical
pressure reaction is exactly the $l=1$ moment $\mathscr B_1$ (same projection as in GA),
so \eqref{eq:linelift} is subsumed into $\mathscr B_1$ rather than added separately, and
$\dot v=-\Bo+\mathscr B_1$ holds unchanged.
\end{proof}

\noindent\emph{Consistency with GA.} For $\theta_d=\pi$ (tangency) $\sin\theta_d=0$ and
$F_{\text{cl}}^z=0$: GA's smooth-cap pressure carries no edge load, as expected. For
finite $\theta_d$ the reaction includes a near-concentrated line load at $\theta_c$.

\begin{remark}[Numerical edge-load caveat]\label{rem:edgeload}
Because the finite-angle reaction \eqref{eq:linelift} concentrates near $\theta_c$, the
truncated degree-$L$ Legendre pressure must represent a near-singular edge load. GA never
faced this ($\sin\theta_d\equiv0$). Expect Gibbs-type oscillation and $O(1/L)$ error in
$\mathscr B_1$ for strongly non-tangential angles; the effect is $O(\sin\theta_d)$ and
hence \emph{small in the near-non-wetting regime} ($\theta_e\to\pi$) that is our
experimental target. The energy-budget closure test (\S\ref{sec:budget}) is the
quantitative bound on this error and should be monitored as $\theta_e$ departs from
$\pi$.
\end{remark}

\noindent\emph{Remark (mechanism).} The wetting energy \eqref{eq:Ewet} changes the
dynamics \emph{indirectly}: it alters the contact-line boundary condition
(\S\ref{sec:mobility}), which changes the admissible shape and hence the pressure
multipliers $\mathscr B_l$ in \eqref{eq:mode} and \eqref{eq:vdot}. The extension is thus
a change of \emph{constraint}, not an added body force --- the key to preserving GA's
linear structure.

\section{The contact-line law}\label{sec:mobility}

Treat the contact-line radius $r_c$ (equivalently $\theta_c$) as a genuine dynamical
variable. Its motion dissipates energy through microscopic and viscous processes near
the moving line whose net effect, at leading order and for small capillary number, is a
force linear in the line speed \citep{Blake2006,deGennes1985,Cox1986,Voinov1976}; we
encode this by a Rayleigh dissipation functional quadratic in the line speed,
\be\label{eq:rayleigh}
\mathcal R=\tfrac12\,\xi\,(2\pi r_c)\,\dot r_c^{\,2}\;\ge0,
\ee
with $\xi\ge0$ the contact-line friction coefficient (units of $\sigma t_\sigma/R$ in
dimensionless form; $\xi\to0$ = perfectly mobile, $\xi\to\infty$ = pinned). The
Lagrange--Rayleigh equation for the massless coordinate $r_c$ is
$\partial\mathcal R/\partial\dot r_c=-\,\partial E/\partial r_c$. We use the geometric
identity
\be\label{eq:dAlg}
\frac{\partial A_{lg}}{\partial A_w}=+\cos\theta_d,
\ee
i.e.\ advancing the wetted area by $\mathrm dA_w$ changes the liquid--gas area by
$+\cos\theta_d\,\mathrm dA_w$; the sign is fixed by requiring that
$\partial E/\partial A_w=0$ at $\theta_d=\theta_e$ reproduce Young's relation
(checks: complete wetting $\theta_d\to0$ gives $+1$; perfect non-wetting
$\theta_d\to\pi$ gives $-1$). With $A_w=\pi r_c^2$ and \eqref{eq:Ewet},
\be
\frac{\partial E}{\partial r_c}
=\frac{\partial}{\partial r_c}\big(A_{lg}-\cos\theta_e\,\pi r_c^2\big)
=2\pi r_c\big(\cos\theta_d-\cos\theta_e\big),
\ee
so that
$\partial\mathcal R/\partial\dot r_c=\xi\,2\pi r_c\,\dot r_c
=-\,\partial E/\partial r_c=2\pi r_c(\cos\theta_e-\cos\theta_d)$, giving the linear
contact-line (Blake/de~Gennes) law
\be\label{eq:mobility}
\boxed{\;\xi\,\dot r_c=\cos\theta_e-\cos\theta_d.\;}
\ee
The right-hand side is the uncompensated Young force per unit length. Equation
\eqref{eq:mobility} \emph{replaces} the tangency condition \eqref{eq:tangency} as the
closure for the contact line.

\begin{corollary}[Quasi-static limit]\label{cor:quasistatic}
As $\xi\to0$, \eqref{eq:mobility} forces $\cos\theta_d=\cos\theta_e$, i.e.\
$\theta_d=\theta_e$: a perfectly mobile line instantaneously adopts the equilibrium
angle.
\end{corollary}
\begin{corollary}[Non-wetting limit]\label{cor:nonwetting}
Setting $\theta_e=\pi$ in the quasi-static limit gives $\theta_d=\pi$, which by
Proposition~\ref{prop:tangency} is exactly the GA tangency condition
\eqref{eq:tangency}. Hence $(\theta_e\to\pi,\ \xi\to0)$ recovers GA.
\end{corollary}
\begin{corollary}[Pinned limit]\label{cor:pinned}
As $\xi\to\infty$, \eqref{eq:mobility} gives $\dot r_c\to0$: the contact line is pinned.
\end{corollary}

\subsection{Advancing, receding, and contact-angle hysteresis}\label{sec:hysteresis}

As written, \eqref{eq:mobility} uses a single friction $\xi$ and a single equilibrium
angle $\theta_e$, so it is odd under $\dot r_c\to-\dot r_c$ about $\theta_e$: the
\emph{advancing} contact line of the spreading phase and the \emph{receding} line of the
retraction phase obey the \emph{same} relation, and the model has no contact-angle
hysteresis. Real substrates instead exhibit distinct advancing and receding angles
$\theta_a\ge\theta_r$ with a pinned band in between, an important and often dominant
source of dissipation and of the spreading/retraction asymmetry seen in rebound
\citep{JoannyDeGennes1984,SnoeijerAndreotti2013}. The present framework accommodates
this at no structural cost by promoting \eqref{eq:mobility} to a piecewise (direction- and
angle-dependent) law,
\be\label{eq:hysteresis}
\xi\,\dot r_c=
\begin{cases}
\cos\theta_a-\cos\theta_d, & \theta_d<\theta_a \quad(\text{advancing},\ \dot r_c>0),\\[2pt]
0, & \theta_r\le\theta_d\le\theta_a \quad(\text{pinned}),\\[2pt]
\cos\theta_r-\cos\theta_d, & \theta_d>\theta_r \quad(\text{receding},\ \dot r_c<0),
\end{cases}
\ee
optionally with distinct $\xi_a,\xi_r$. Equation \eqref{eq:mobility} is the
zero-hysteresis case $\theta_a=\theta_r=\theta_e$. The single-angle law is retained as the
baseline here because it isolates the leading-order effect of wettability with one fitted
parameter; \eqref{eq:hysteresis} is the natural refinement once advancing and receding
angles are measured independently, and both branches are represented in the same discrete
selection residual (\S\ref{sec:discrete}) by the sign of the contact-line increment.

\section{Energy budget}\label{sec:budget}

\begin{proposition}[Dissipation balance]\label{prop:budget}
Along solutions of \eqref{eq:kin}--\eqref{eq:vdot} closed by \eqref{eq:mobility},
\be\label{eq:budget}
\frac{\mathrm dE}{\mathrm dt}
=-\underbrace{\sum_{l\ge2}2(2l+1)(l-1)\,\Oh\,\mathscr U_l^2}_{\text{bulk viscous }\ge0}
 \;-\;\underbrace{\xi\,(2\pi r_c)\,\dot r_c^{\,2}}_{\text{contact line }\ge0}\;\le0,
\ee
with $E$ from \eqref{eq:energy}. In the inviscid, perfectly-mobile, non-wetting limit
($\Oh\to0,\ \xi\to0,\ \theta_e\to\pi$) energy is conserved, reproducing the GA
inertio-capillary limit.
\end{proposition}
\begin{proof}
The pressure does no net work on the frictionless unilateral constraint, leaving the
bulk viscous sink; the wetting and moving-line liquid--gas terms combine, via the
mobility law \eqref{eq:mobility}, into the contact-line sink. The full calculation is
given in Appendix~\ref{app:budget}.
\end{proof}

Equation \eqref{eq:budget} is the arbiter used to validate the implementation: it must
close numerically, and it certifies that no spurious force (e.g.\ a vertical Young term
in \eqref{eq:vdot}) has been introduced --- consistent with
Proposition~\ref{prop:noforce}.

\section{Discrete formulation}\label{sec:discrete}

We inherit GA Appendix~B verbatim for the bulk: truncation at mode $L$, implicit-Euler
(or BDF) time stepping, the angular mesh $\{\theta_0=0\}\cup\{\text{zeros of }P_L\}$, and
the block linear system $\mathbf M^{k+1,q}\,\mathbf x^{k+1,q}=\mathbf b^{k+1,q}$ for a
candidate number $q$ of contact mesh points (GA eqns for $\mathbf M,\mathbf x,\mathbf
b$). \emph{The extension leaves $\mathbf M$ and $\mathbf b$ unchanged.} Only the
selection of $q$ and, dynamically, its allowed rate of change are modified.

\subsection{Quasi-static finite angle (perfectly mobile, $\xi=0$)}\label{sec:disc_qs}
GA select $q$ by minimising the tangency residual (their $e^2_q$, the discrete
$\partial_\theta(\zeta\cos\theta)$ across the contact edge, i.e.\ $\propto -z_\theta$).
For a finite target angle $\theta_e$ the condition $\theta_d=\theta_e$ is, by
\eqref{eq:thetad}, that the tangent $(x_\theta,z_\theta)$ be parallel to
$(\cos(\pi-\theta_e),\sin(\pi-\theta_e))=(-\cos\theta_e,\sin\theta_e)$; taking the
(vanishing) cross product gives the \emph{linear, well-conditioned} residual
\be\label{eq:disc_angle}
e^2_q(\theta_e)=\big|\,x_\theta\,\sin\theta_e+z_\theta\,\cos\theta_e\,\big|
\Big|_{\theta=(\theta_q+\theta_{q+1})/2},
\ee
with $x_\theta,z_\theta$ from \eqref{eq:tangent} evaluated on the solved amplitudes
$A_l^{k+1,q}$. This form is preferable to $|\theta_d^q-\theta_e|$ on three counts: it is
\emph{linear} in the amplitudes (for fixed $\theta_q$), it is \emph{bounded} (no division
by $x_\theta$, so it survives the $x_\theta\to0$ spreading regime of
Remark~\ref{rem:xtheta}), and at $\theta_e=\pi$ it collapses \emph{term for term} to GA's
tangency residual: $\sin\theta_e=0,\ \cos\theta_e=-1$ give
$e^2_q=|z_\theta|=|\partial_\theta(\zeta\cos\theta)|$, the discrete finite difference GA
already use. Impenetrability $e^1_q$ (GA) is retained unchanged;
$e_q=\max\{e^1_q,e^2_q\}$ and $m^{k+1}=\operatorname{argmin}_q e_q$.

Thus for $\theta_e=\pi$ the scheme shares GA's tangency root and reduces to GA's discrete
objective; for finite $\theta_e$ only the scalar selection residual changes. No matrix
$\mathbf M$ or right-hand-side $\mathbf b$ change is required, so the linear solve --- and
its conditioning --- is that of GA, subject to the edge-load caveat of
Remark~\ref{rem:edgeload} at angles far from $\pi$.

\subsection{Dynamic contact line (finite $\xi$)}\label{sec:disc_dynamic}
A finite friction must be represented \emph{sub-mesh}: on the fixed angular grid the
contact count $q$ is an integer, so the naive choice of scoring each $q$ by the discrete
mobility imbalance pins the line for any $\xi>0$ (a one-node change $\delta r_c$ makes
$\xi\,\delta r_c$ dominate the $\mathcal O(\delta_t)$ driving term). We instead carry the
contact-line radius $r_c$ as a \emph{continuous} tracker advanced each step by a
rate-limited integration of \eqref{eq:mobility} toward the quasi-static target
$r_c^{\text{qs}}$ (the edge radius selected by \eqref{eq:disc_angle}):
\be\label{eq:disc_mobility}
r_c^{k+1}=r_c^{k}+\operatorname{clamp}\!\Big(\tfrac{\delta_t^{k}}{\xi}\,(\cos\theta_e-\cos\theta_d^{k}),\;
0,\;r_c^{\text{qs}}-r_c^{k}\Big),
\ee
with $\theta_d^{k}$ the apparent angle at the current line; the integer contact set for
the bulk solve is then the mesh bracket of $r_c^{k+1}$. The clamp gives the exact limits:
$\xi\to0$ reaches $r_c^{\text{qs}}$ in one step (quasi-static, \S\ref{sec:disc_qs});
$\xi\to\infty$ leaves $r_c$ unchanged (pinned). The bulk solve at fixed bracket remains
the GA linear system --- the only new work is the scalar tracker update, so the
well-conditioned inner solve responsible for GA's robustness is untouched, and no energy
is injected into the ballistic centre-of-mass motion (verified below via the coefficient
of restitution).

\medskip\noindent\emph{Optional monolithic form.} Alternatively $\theta_c$ may be
carried as a $(3L{+}2)$-th unknown with \eqref{eq:mobility} as an added residual row,
solved by Newton with the analytic bulk Jacobian augmented by one finite-differenced
column. The operator-split tracker above is the default because it cannot destabilise the
proven linear core and, being sub-mesh, already resolves the continuous contact-line
motion the monolithic form would provide.

\section{Numerical validation}\label{sec:validation}

We verify the extended solver against its own limits and check the predicted trends.
Except where noted, results use $L=20$ modes, $\Oh=0.1$, $\Bo=0.01$ and impact speed
$\We^{1/2}=0.5$.

\emph{(i) Exact reduction.} With $(\theta_e=\pi,\ \xi=0)$ the extended solver reproduces
the GA trajectory to machine precision (the selection residual \eqref{eq:disc_angle}
collapses term-for-term to GA's, Appendix~\ref{app:discrete}), confirmed as a
regression test.

\emph{(ii) Wettability lowers restitution.} Decreasing $\theta_e$ from $\pi$ at $\xi=0$
monotonically decreases the coefficient of restitution and increases the contact time
and maximal contact radius:
\begin{center}
\begin{tabular}{lcccc}
$\theta_e/\pi$ & $1.00$ & $0.96$ & $0.92$ & $0.90$\\[2pt]
COR            & $0.657$ & $0.654$ & $0.552$ & $0.361$\\
$t_c$          & $2.60$ & $2.63$ & $2.96$ & $3.40$\\
$\beta_c$      & $0.516$ & $0.516$ & $0.654$ & $0.683$
\end{tabular}
\end{center}
consistent with the loss of rebound energy to increased spreading as the surface becomes
more wettable.

\emph{(iii) Friction resists capture.} At fixed $\theta_e=0.9\pi$, increasing the
contact-line friction $\xi$ monotonically increases the coefficient of restitution,
from the quasi-static value ($0.361$ at $\xi=0$) up to the perfectly non-wetting value
($0.657$) once the line is effectively pinned ($\xi\gtrsim1$); the maximal contact
radius correspondingly falls back to its non-wetting value. This realises the
$\xi\to0$ (quasi-static) and $\xi\to\infty$ (pinned) limits of \eqref{eq:disc_mobility}
and shows contact-line friction acting as a distinct, tunable dissipation channel.

\emph{(iv) No energy injection.} In every case above the coefficient of restitution
satisfies $\text{COR}\le1$, the operational form of the energy budget
(Prop.~\ref{prop:budget}): no spurious energy enters the ballistic centre-of-mass
motion, including at the finite contact angles where the edge-load truncation error
\eqref{eq:linelift} is largest. Quantitative accuracy is expected to be best near
$\theta_e=\pi$ (error $\mathcal O(\sin\theta_d)$, Remark~\ref{rem:edgeload}); the
trustworthy band for the parameters here is $\theta_e\gtrsim0.9\pi$.

\section{Summary of the extended model}\label{sec:summary}

The extended KM model is \eqref{eq:kin}--\eqref{eq:vdot} together with the contact
constraints \eqref{eq:contact}--\eqref{eq:pzero} and the contact-line law
\eqref{eq:mobility} in place of the tangency condition \eqref{eq:tangency}, with
$\theta_d$ given by \eqref{eq:thetad}. It has:
\begin{itemize}
\item one new dynamical degree of freedom, the contact-line radius $r_c$;
\item two physical parameters, the equilibrium angle $\theta_e$ (directly measurable)
      and the contact-line friction $\xi$ (a single fit parameter);
\item an exact reduction to GA at $(\theta_e\to\pi,\ \xi\to0)$
      (Cor.~\ref{cor:nonwetting});
\item a COM momentum equation that retains the GA form, the contact-line reaction being
      absorbed into $\mathscr B_1$ (Prop.~\ref{prop:noforce});
\item a closed energy budget with a distinct contact-line dissipation channel
      (Prop.~\ref{prop:budget});
\item a discrete form that changes only the contact-set selection criterion
      (\S\ref{sec:discrete}), preserving GA's linear solve.
\end{itemize}

\appendix

\section{Full derivation of the energy budget}\label{app:budget}

We give the calculation sketched in Proposition~\ref{prop:budget}. Differentiate the
kinetic and capillary contributions of \eqref{eq:energy}. Using the modal energy
$e_l=\tfrac12\mathscr U_l^2+\tfrac12 l(l+2)(l-1)\mathscr A_l^2$ and the equations of
motion \eqref{eq:kin}--\eqref{eq:mode},
\be
\dot e_l=\mathscr U_l\dot{\mathscr U}_l+l(l+2)(l-1)\mathscr A_l\dot{\mathscr A}_l
=-2(2l+1)(l-1)\,\Oh\,\mathscr U_l^2-\,l\,\mathscr B_l\mathscr U_l,
\ee
so summing over $l\ge2$ and adding $\dot E_{\text{grav}}+\tfrac12\frac{\mathrm d}{\mathrm
dt}v^2=\Bo\,v+v\dot v=v\mathscr B_1$ (from \eqref{eq:vdot}) gives
\be\label{eq:bulkwork}
\dot E_{\text{bulk}}
=-\sum_{l\ge2}2(2l+1)(l-1)\,\Oh\,\mathscr U_l^2
-\Big(\sum_{l\ge2}l\,\mathscr B_l\mathscr U_l-v\,\mathscr B_1\Big).
\ee
The parenthesised term is the rate of working of the contact pressure on the interface.
On the contact set the constraint fixes the normal position ($z=0$), so the interface
normal velocity vanishes wherever $p\neq0$; off the contact set $p=0$. Hence
$\int p\,u_n\,\mathrm dS=0$ and the pressure-work term vanishes identically (frictionless
unilateral contact does no net work), leaving only the bulk viscous sink. Finally the
surface--wetting energy: from \eqref{eq:Ewet} and \eqref{eq:dAlg},
\be
\dot E_{\text{surf}}
=\underbrace{\frac{\partial A_{lg}}{\partial r_c}\dot r_c}_{+2\pi r_c\cos\theta_d\,\dot r_c}
-\cos\theta_e\,2\pi r_c\,\dot r_c
=-2\pi r_c(\cos\theta_e-\cos\theta_d)\,\dot r_c
=-\xi(2\pi r_c)\,\dot r_c^2,
\ee
the last equality by the mobility law \eqref{eq:mobility}. Adding the two contributions
yields \eqref{eq:budget}. Both sinks are manifestly non-negative, and each vanishes in
its respective limit ($\Oh\to0$; $\xi\to0$ or $\dot r_c\to0$).

\section{The discrete selection residual and its reduction to GA}\label{app:discrete}

GA (their Appendix~B) solve, for each candidate contact count $q$, the linear system
$\mathbf M^{k+1,q}\mathbf x^{k+1,q}=\mathbf b^{k+1,q}$ and select $q$ by minimising the
tangency residual $e^2_q=\big|\,Z(\theta_q)-Z(\theta_{q+1})\,\big|$, with
$Z(\theta)=\zeta(\theta)\cos\theta+z$ the collocated height. Our extension leaves
$\mathbf M,\mathbf b$ untouched and replaces only $e^2_q$. Writing the discrete edge
differences $\Delta Z=Z(\theta_q)-Z(\theta_{q+1})$ and
$\Delta x=x(\theta_q)-x(\theta_{q+1})$ with $x(\theta)=\zeta(\theta)\sin\theta$, the
finite-angle residual is
\be\label{eq:disc_residual_app}
e^2_q(\theta_e)=\big|\,\Delta x\,\sin\theta_e+\Delta Z\,\cos\theta_e\,\big|,
\ee
the discrete form of \eqref{eq:disc_angle}. Two properties follow directly. (i)
\emph{Exact GA reduction:} at $\theta_e=\pi$, $\sin\theta_e=0$ and $\cos\theta_e=-1$, so
$e^2_q(\pi)=|\Delta Z|$, GA's tangency residual term for term; the selected $q$, and hence
the entire trajectory, are identical. (ii) \emph{Boundedness:} \eqref{eq:disc_residual_app}
is a linear combination of collocated amplitude differences, with no division by $x_\theta$,
so it remains well defined through the maximal-spreading regime where $x_\theta\to0$
(Remark~\ref{rem:xtheta}). For the hysteretic law \eqref{eq:hysteresis}, the same residual
is used with $\theta_e$ replaced by $\theta_a$ or $\theta_r$ according to the sign of the
trial contact-line increment $r_c^{k+1,q}-r_c^{k}$, and pinned candidates
($\theta_r\le\theta_d\le\theta_a$) are scored by their distance to the nearest branch.

\bibliographystyle{jfm}
\begin{thebibliography}{}

\bibitem[Alventosa \emph{et~al.}(2023)]{Alventosa2023}
Alventosa, L.F.L., Cimpeanu, R. \& Harris, D.M. 2023 Inertio-capillary rebound of a
droplet impacting a fluid bath. \emph{J. Fluid Mech.} \textbf{958}, A24.

\bibitem[Blake(2006)]{Blake2006}
Blake, T.D. 2006 The physics of moving wetting lines. \emph{J. Colloid Interface Sci.}
\textbf{299}, 1--13.

\bibitem[Cox(1986)]{Cox1986}
Cox, R.G. 1986 The dynamics of the spreading of liquids on a solid surface. Part 1.
Viscous flow. \emph{J. Fluid Mech.} \textbf{168}, 169--194.

\bibitem[Gabbard \emph{et~al.}(2025)]{gabbard2025}
Gabbard, C.T., Ag\"uero, E.A., Cimpeanu, R., Kuehr, K., Silver, E., Barotta, J.-W.,
Galeano-Rios, C.A. \& Harris, D.M. 2025 Drop rebound at low Weber number.
\emph{J. Fluid Mech.} (companion paper).

\bibitem[Galeano-Rios \emph{et~al.}(2017)]{GaleanoRios2017}
Galeano-Rios, C.A., Milewski, P.A. \& Vanden-Broeck, J.-M. 2017 Non-wetting impact of a
sphere onto a bath and its application to bouncing droplets. \emph{J. Fluid Mech.}
\textbf{826}, 97--127.

\bibitem[de Gennes(1985)]{deGennes1985}
de Gennes, P.-G. 1985 Wetting: statics and dynamics. \emph{Rev. Mod. Phys.} \textbf{57},
827--863.

\bibitem[Grmela \& \"Ottinger(1997)]{GrmelaOttinger1997}
Grmela, M. \& \"Ottinger, H.C. 1997 Dynamics and thermodynamics of complex fluids. I.
Development of a general formalism. \emph{Phys. Rev. E} \textbf{56}, 6620--6632.

\bibitem[Joanny \& de Gennes(1984)]{JoannyDeGennes1984}
Joanny, J.F. \& de Gennes, P.-G. 1984 A model for contact angle hysteresis.
\emph{J. Chem. Phys.} \textbf{81}, 552--562.

\bibitem[Lamb(1932)]{Lamb1924}
Lamb, H. 1932 \emph{Hydrodynamics}, 6th edn. Cambridge University Press.

\bibitem[Mol\'a\v{c}ek \& Bush(2013)]{MolacekBush2013}
Mol\'a\v{c}ek, J. \& Bush, J.W.M. 2013 Drops bouncing on a vibrating bath.
\emph{J. Fluid Mech.} \textbf{727}, 582--611.

\bibitem[Okumura \emph{et~al.}(2003)]{Okumura2003}
Okumura, K., Chevy, F., Richard, D., Qu\'er\'e, D. \& Clanet, C. 2003 Water spring: a
model for bouncing drops. \emph{Europhys. Lett.} \textbf{62}, 237--243.

\bibitem[Richard \emph{et~al.}(2002)]{Richard2002}
Richard, D., Clanet, C. \& Qu\'er\'e, D. 2002 Contact time of a bouncing drop.
\emph{Nature} \textbf{417}, 811.

\bibitem[Snoeijer \& Andreotti(2013)]{SnoeijerAndreotti2013}
Snoeijer, J.H. \& Andreotti, B. 2013 Moving contact lines: scales, regimes, and
dynamical transitions. \emph{Annu. Rev. Fluid Mech.} \textbf{45}, 269--292.

\bibitem[Voinov(1976)]{Voinov1976}
Voinov, O.V. 1976 Hydrodynamics of wetting. \emph{Fluid Dyn.} \textbf{11}, 714--721.

\end{thebibliography}

\end{document}
